\chapter{Vector Fields and Flows}

\begin{center}\itshape
A vector field is an arrow at every point; follow the arrows and you get a flow, a motion of the whole space. Two flows run one after the other in
different orders can end in different places, and the infinitesimal size of the discrepancy is the Lie bracket. When all the brackets vanish, the flows
fit together into coordinates: that is Frobenius' theorem. And on an even-dimensional sphere no flow can move every point, which is the hairy ball
theorem.
\end{center}

In July 1870 a young Norwegian geometer was in Paris, working with his friend Felix Klein, when France declared war on Prussia. Klein, a Prussian, left at
once. Sophus Lie stayed a little longer and then set off on foot towards Italy. In August, near Fontainebleau, he was arrested as a German spy: he was a
foreigner walking alone, and his papers were full of letters and symbols that the police took for a code. He spent about a month in prison, working on
his mathematics, until Gaston Darboux came from Paris and obtained his release. Lie liked to say afterwards that a mathematician is well equipped to
bear prison.

What Lie carried in his notes was the beginning of an idea that would occupy him for the rest of his life. Galois had shown that the solvability of an
algebraic equation is governed by a finite group of symmetries. Lie wanted the same for differential equations, and the symmetries he found were not
finite: they were \emph{continuous}, families of transformations depending on parameters, generated by infinitesimal motions. An infinitesimal motion
is a vector field; the transformations it generates are its flow; and the way two infinitesimal motions fail to commute is the bracket that now bears
his name. Chapter~7 is about the groups; this chapter is about the vector fields and their flows.

We start with the existence and uniqueness of solutions of differential equations, and with their smooth dependence on the initial point, which no
earlier volume proves and everything later needs. Then we move to manifolds, define the bracket in three equivalent ways, and prove Frobenius' theorem.
We end with Milnor's 1978 proof of the hairy ball theorem, which uses nothing but a change of variables and the fact that $(1+t^2)^{3/2}$ is not a
polynomial.

\section{Differential equations in \texorpdfstring{$\R^n$}{R\^{}n}}

Let $U\subset\R^n$ be open and $F\colon U\to\R^n$ a $C^1$ map, which we think of as a field of velocities. A \emph{solution} or \emph{integral curve}
is a $C^1$ map $x\colon J\to U$ on an interval $J$ with $x'(t)=F(x(t))$. Equations whose right side depends on $t$ reduce to this case by adding $t$ as
an extra coordinate with $t'=1$, so we lose nothing by considering only autonomous equations. A $C^1$ map is locally Lipschitz: on a closed ball
$\bar B\subset U$, $\|DF\|$ is bounded by some $L$ and Lemma~\ref{lem:mvi} gives $|F(x)-F(y)|\le L|x-y|$.

\begin{lemma}[Gr\"onwall]\label{lem:gronwall}
Let $u\colon[0,T]\to[0,\infty[$ be continuous, $a,b\ge0$, and $u(t)\le a+b\int_0^tu(s)\dd s$ for all $t$. Then $u(t)\le ae^{bt}$.
\end{lemma}
\begin{proof}
Let $w(t)=a+b\int_0^tu$. Then $w'=bu\le bw$, so $(e^{-bt}w)'\le0$ and $w(t)\le w(0)e^{bt}=ae^{bt}$. Since $u\le w$, we are done.
\end{proof}

\begin{theorem}[Picard--Lindel\"of]\label{thm:picard}
Let $x_0\in U$, $\bar B(x_0,b)\subset U$, $M=\max_{\bar B}|F|$ and $L$ a Lipschitz constant of $F$ on $\bar B(x_0,b)$. For $\tau=\min(\frac bM,\frac1{2L})$
there is a solution $x\colon[-\tau,\tau]\to\bar B(x_0,b)$ with $x(0)=x_0$. Two solutions defined on the same interval containing $0$ and with the same value
at $0$ are equal.
\end{theorem}
\begin{proof}
\emph{Existence.} Let $\mathcal C$ be the set of continuous maps $[-\tau,\tau]\to\bar B(x_0,b)$ with the sup distance; it is complete, as a closed subset of
a Banach space (Volume~7, Proposition~2.7). For $x\in\mathcal C$ let $(Tx)(t)=x_0+\int_0^tF(x(s))\dd s$. Then $|Tx(t)-x_0|\le M|t|\le b$, so $T$ maps
$\mathcal C$ into itself, and
\[
|Tx(t)-Ty(t)|\le\Bigl|\int_0^tL|x(s)-y(s)|\dd s\Bigr|\le L\tau\,d(x,y)\le\tfrac12d(x,y).
\]
By Banach's fixed point theorem $T$ has a fixed point $x$. Since $s\mapsto F(x(s))$ is continuous, $x=Tx$ is $C^1$ with $x'=F(x)$ and $x(0)=x_0$.

\emph{Uniqueness.} Let $x,y\colon J\to U$ be solutions with $x(0)=y(0)$, and $[0,T]\subset J$ (negative times are symmetric). On the compact set
$x([0,T])\cup y([0,T])$ the map $F$ is Lipschitz with some constant $L'$, by the remark that follows. Then $u=|x-y|$ satisfies $u(t)\le\int_0^tL'u$, and Lemma~\ref{lem:gronwall} with $a=0$ gives $u=0$.
\end{proof}

The remark is the following fact, which we use repeatedly: a $C^1$ map is Lipschitz on every compact subset
$K$ of its domain. If not, there would be $x_k,y_k\in K$ with $|F(x_k)-F(y_k)|>k|x_k-y_k|$; since $F$ is bounded on $K$, $|x_k-y_k|\to0$, and after
extracting, both sequences converge to some $z\in K$, so that eventually $x_k,y_k$ lie in a ball around $z$ on which $F$ is Lipschitz, a contradiction.

\begin{think}
The equation $x'=\sqrt{|x|}$ has the solution $x=0$, but also $x(t)=\frac{t^2}4$ for $t\ge0$ (and $0$ before). Which hypothesis of
Theorem~\ref{thm:picard} fails? Can you find infinitely many solutions with $x(0)=0$?
\end{think}

By uniqueness, all solutions through $x_0$ at time $0$ agree where they are both defined, so their union is a solution on an open interval
$J(x_0)\ni0$, the \emph{maximal solution} $t\mapsto\varphi(t,x_0)$. The \emph{flow} of $F$ is the map $\varphi$ on the \emph{flow domain}
$\mathcal D=\{(t,x):x\in U,\ t\in J(x)\}$.

\begin{proposition}[escape lemma]\label{prop:escape}
If $J(x_0)$ has a finite upper end $\beta$, then the maximal solution leaves every compact $K\subset U$: there is no compact $K$ with
$\varphi(t,x_0)\in K$ for all $t\in[0,\beta[$.
\end{proposition}
\begin{proof}
Let $K$ be compact, choose $b>0$ with $K_b=\{x:d(x,K)\le b\}\subset U$, and let $M$ and $L$ be bounds for $|F|$ and a Lipschitz constant on the compact
$K_b$. For every $y\in K$, $\bar B(y,b)\subset K_b$, so Theorem~\ref{thm:picard} gives a solution starting at $y$ on $[-\tau,\tau]$ with $\tau$ independent of
$y$. If the maximal solution stayed in $K$, the solution starting at $\varphi(\beta-\frac\tau2,x_0)$ would extend it up to time $\beta+\frac\tau2$,
contradicting maximality.
\end{proof}

\begin{theorem}[continuous dependence]\label{thm:flowcont}
The flow domain $\mathcal D$ is open in $\R\times U$ and the flow $\varphi$ is continuous on it. More precisely, if $[0,T]\subset J(x_0)$, there are
$\delta>0$ and $L$ such that for $|x-x_0|<\delta$ we have $[0,T]\subset J(x)$ and $|\varphi(t,x)-\varphi(t,x_0)|\le e^{Lt}|x-x_0|$ for $t\in[0,T]$.
\end{theorem}
\begin{proof}
Let $K=\varphi([0,T],x_0)$, compact, and choose $b>0$ with $K_b\subset U$ and $L$ a Lipschitz constant for $F$ on $K_b$. Take
$\delta=\frac b2e^{-LT}$ and $|x-x_0|<\delta$, and let $t^*$ be the supremum of the $t\in[0,T]\cap J(x)$ with $\varphi([0,t],x)\subset K_b$. For $s<t^*$, the
function $u(s)=|\varphi(s,x)-\varphi(s,x_0)|$ satisfies $u(s)\le|x-x_0|+\int_0^sLu$, so $u(s)\le e^{Ls}|x-x_0|<\frac b2$ by Lemma~\ref{lem:gronwall}: the solution
stays in the compact $K_{b/2}$ on $[0,t^*[$. By Proposition~\ref{prop:escape}, $t^*$ is not the end of $J(x)$, so $t^*\in J(x)$ and $\varphi(t^*,x)\in K_{b/2}$, which lies
in the interior of $K_b$. If $t^*<T$, the solution would stay in $K_b$ a little after $t^*$, contradicting the definition of $t^*$. So $t^*=T$, $[0,T]\subset J(x)$
and the estimate holds on $[0,T]$. Negative times are symmetric.

Now let $(t_0,x_0)\in\mathcal D$, say $t_0\ge0$, and choose $T>t_0$ and $-\varepsilon<0$ in the open interval $J(x_0)$. The estimate on $[0,T]$ and its
symmetric version on $[-\varepsilon,0]$ show that $[-\varepsilon,T]\times B(x_0,\delta)\subset\mathcal D$ for a small $\delta$, so $\mathcal D$ is open. Finally
$|\varphi(t,x)-\varphi(t_0,x_0)|\le e^{LT}|x-x_0|+|\varphi(t,x_0)-\varphi(t_0,x_0)|$, which tends to $0$.
\end{proof}

\subsection*{Smooth dependence on the initial point}

How does the endpoint $\varphi(t,x)$ change when we move the starting point $x$? To first order, a small displacement $h$ of the start is carried along by
the linearized equation.

\begin{theorem}[smooth dependence]\label{thm:flowsmooth}
If $F$ is $C^k$, $1\le k\le\infty$, then the flow $\varphi$ is $C^k$ on $\mathcal D$. Its derivative $Y(t)=D_x\varphi(t,x)$ is the solution of the
\emph{variational equation}
\[
Y'(t)=DF(\varphi(t,x))\,Y(t),\qquad Y(0)=I.
\]
\end{theorem}
\begin{proof}
\emph{The candidate.} Fix $x_0$ and $[0,T]\subset J(x_0)$, and take $\delta$, $b$, $L$ as in Theorem~\ref{thm:flowcont}. The variational equation is linear
with continuous coefficients, so it has a solution $Y_{x}$ on $[0,T]$ for each $|x-x_0|<\delta$ (Theorem~\ref{thm:picard} applied to a linear equation
gives solutions on intervals of fixed length, which we can concatenate).

\emph{Differentiability.} The derivative $DF$ is uniformly continuous on the compact $K_b$, so there is a function $\omega$ with $\omega(r)\to0$ as $r\to0$
and $|F(a+u)-F(a)-DF(a)u|\le\omega(|u|)\,|u|$ when $a$ and $a+u$ lie in $K_b$ and the segment between them too (Lemma~\ref{lem:mvi} applied to
$z\mapsto F(z)-DF(a)z$). For small $h$, put $u(t)=\varphi(t,x_0+h)-\varphi(t,x_0)$, so $|u(t)|\le e^{LT}|h|$, and $r(t)=u(t)-Y_{x_0}(t)h$. Then $r(0)=0$ and
\[
r'(t)=DF(\varphi(t,x_0))\,r(t)+\rho(t),\qquad|\rho(t)|\le\omega(e^{LT}|h|)\,e^{LT}|h|.
\]
With $M=\max_{K_b}\|DF\|$ we get $|r(t)|\le\int_0^t(M|r|+|\rho|)\le T\max|\rho|+M\int_0^t|r|$, and Lemma~\ref{lem:gronwall} gives
$|r(t)|\le Te^{MT}e^{LT}\omega(e^{LT}|h|)\,|h|=o(|h|)$. So $D_x\varphi(t,x_0)=Y_{x_0}(t)$.

\emph{Continuity of the derivative.} Let $A_x(t)=DF(\varphi(t,x))$. As $x\to x_0$, $\varphi(t,x)\to\varphi(t,x_0)$ uniformly in $t\in[0,T]$ by
Theorem~\ref{thm:flowcont}, so $\eta(x)=\sup_t\|A_x(t)-A_{x_0}(t)\|\to0$ by uniform continuity of $DF$ on $K_b$. The difference $Z=Y_x-Y_{x_0}$ satisfies
$Z'=A_xZ+(A_x-A_{x_0})Y_{x_0}$, $Z(0)=0$, and Gr\"onwall gives $\sup_t\|Z(t)\|\le T\eta(x)\max\|Y_{x_0}\|e^{MT}\to0$. Together with the continuity of each
$Y_x$ in $t$, this makes $(t,x)\mapsto D_x\varphi(t,x)$ continuous. Since $\partial_t\varphi=F\circ\varphi$ is continuous too, all partial derivatives of
$\varphi$ are continuous and $\varphi$ is $C^1$ (Proposition~\ref{prop:partials}).

\emph{Higher derivatives.} We prove by induction on $k$ that the flow of every $C^k$ field is $C^k$; the case $k=1$ is done. Let $F$ be $C^k$ with
$k\ge2$. The field $G(x,Y)=(F(x),DF(x)Y)$ on $U\times M_n(\R)$ is $C^{k-1}$, and by what we proved its flow is
$(t,x,Y_0)\mapsto(\varphi(t,x),D_x\varphi(t,x)Y_0)$. By the induction hypothesis this is $C^{k-1}$, so $D_x\varphi$ is $C^{k-1}$ (take $Y_0=I$). Also
$\varphi$ is $C^{k-1}$, being the flow of the $C^{k-1}$ field $F$, so $\partial_t\varphi=F\circ\varphi$ is $C^{k-1}$. All first partial derivatives of $\varphi$
are $C^{k-1}$, so $\varphi$ is $C^k$.
\end{proof}

The variational equation is also how one computes. It says how an error in the initial data grows: a perturbation $h$ becomes $Y(t)h$, and the norm of
$Y(t)$ is the factor by which errors are magnified. When it grows exponentially the system is chaotic; for the simple pendulum it grows only
linearly.

\section{Vector fields on manifolds and their flows}

\begin{definition}
A \emph{vector field} on a smooth manifold $M$ is a choice of $X_p\in T_pM$ for every $p$ which is smooth: in every chart,
$X=\sum_iX^i\frac{\partial}{\partial x^i}$ with smooth coefficients $X^i$. It acts on functions by $(Xf)(p)=X_p(f)$, which is again smooth. An
\emph{integral curve} is a smooth $\gamma\colon J\to M$ with $\gamma'(t)=X_{\gamma(t)}$, where $\gamma'(t)=d\gamma_t(\frac{\dd}{\dd t})$.
\end{definition}

In a chart, the equation $\gamma'=X_\gamma$ reads $(x^i\circ\gamma)'=X^i(\gamma)$, an equation of the previous section. So integral curves exist locally
and are unique locally. Uniqueness is global, by the Hausdorff property: if two integral curves on an interval $J$ agree at one time, the set of times where
they agree is non-empty, closed (the diagonal of $M\times M$ is closed in a Hausdorff space) and open (local uniqueness), so it is $J$. Thus each $p$ has a
maximal integral curve $t\mapsto\varphi_t(p)$ defined on an open interval $J(p)$, and the flow property
\[
\varphi_t(\varphi_s(p))=\varphi_{t+s}(p)
\]
holds whenever the right side is defined, because both sides are integral curves in $t$ with the same value at $t=0$. A vector field is \emph{complete} if
every $J(p)$ is $\R$.

\begin{theorem}\label{thm:complete}
On a compact manifold, every vector field is complete, and its flow $(t,p)\mapsto\varphi_t(p)$ is smooth on $\R\times M$. Each $\varphi_t$ is a diffeomorphism
with inverse $\varphi_{-t}$, and $t\mapsto\varphi_t$ is a homomorphism from $(\R,+)$ to the group of diffeomorphisms of $M$. The same holds on any
manifold for a vector field that vanishes outside a compact set.
\end{theorem}
\begin{proof}
\emph{A uniform time.} For each $p\in M$ take a chart $\kappa$ around $p$. By Theorems~\ref{thm:flowcont} and~\ref{thm:flowsmooth} applied in the chart, the flow
domain of the chart equation is open and contains $(0,\kappa(p))$, so it contains $]-\varepsilon_p,\varepsilon_p[\times B$ for some ball $B$ around $\kappa(p)$,
on which the flow is smooth. Let $V_p$ be the preimage of $B$. Finitely many $V_p$ cover the compact $M$ (or the compact support of $X$; outside it, the
integral curves are constant). Let $\varepsilon$ be the smallest of the corresponding $\varepsilon_p$. Then every integral curve exists at least on
$]-\varepsilon,\varepsilon[$, and $(t,p)\mapsto\varphi_t(p)$ is smooth on $]-\varepsilon,\varepsilon[\times M$, since near each point it coincides with a chart flow.

\emph{All times.} If $J(p)$ had a finite upper end $\beta$, the integral curve from $\varphi_{\beta-\varepsilon/2}(p)$ would extend the maximal one to
$\beta+\frac\varepsilon2$. So $J(p)=\R$. For $(t_0,p_0)$ fixed, write $t_0=N\frac{t_0}N$ with $|\frac{t_0}N|<\varepsilon$; then $\varphi_{t_0}=(\varphi_{t_0/N})^N$ is
smooth, and on $]t_0-\varepsilon,t_0+\varepsilon[\times M$ we have $\varphi(t,p)=\varphi(t-t_0,\varphi_{t_0}(p))$, a composite of smooth maps. Finally
$\varphi_t\circ\varphi_{-t}=\varphi_0=\mathrm{id}$.
\end{proof}

\begin{remark}
On a general manifold the flow domain is an open subset of $\R\times M$ and the flow is smooth on it; the proof splits a compact piece of an integral
curve into short arcs, each inside a chart, and composes the chart flows (Lee, Theorem~9.12). We will only need complete flows.
\end{remark}

\begin{example}[linear flows on the torus]\label{ex:torusflow}
On $T^2=\R^2/\Z^2$ the constant field $X=\frac\partial{\partial\theta_1}+\alpha\frac\partial{\partial\theta_2}$ has the flow
$\varphi_t(\theta)=\theta+t(1,\alpha)$. If $\alpha=\frac pq$ is rational, every orbit closes after time $q$. If $\alpha$ is irrational, every orbit is dense:
the times $t=k\in\Z$ visit the points $(\theta_1,\theta_2+k\alpha)$, and $k\alpha$ modulo $1$ is dense (it is even equidistributed, by Weyl's criterion, Volume~6, Chapter~11). Figure~\ref{fig:torusflow}
shows an orbit for the golden slope. The orbit is the image of an injective immersion $\R\to T^2$ which is not a submanifold, a cousin of the
figure eight of Chapter~1.
\end{example}

\begin{figure}[ht]
\centering
\fig{ch02-torusflow}
\caption{An orbit of the constant field $(1,\alpha)$ on the torus, with $\alpha=\frac{\sqrt5-1}2$, drawn in the unit square whose opposite sides are glued
(as the arrows indicate). The orbit starts at the rust dot and darkens with time. It never closes, and in the long run it fills the square uniformly.}
\label{fig:torusflow}
\end{figure}

\section{The Lie bracket}

A vector field differentiates functions, so two of them can be composed: $f\mapsto X(Yf)$. This is not a vector field, since it involves second
derivatives, but the second derivatives cancel in the commutator.

\begin{proposition}\label{prop:bracket}
For vector fields $X,Y$ on $M$, the operator $[X,Y]f=X(Yf)-Y(Xf)$ is a vector field. In a chart,
\[
[X,Y]^i=\sum_j\bigl(X^j\partial_jY^i-Y^j\partial_jX^i\bigr),
\]
that is, $[X,Y]=DY\cdot X-DX\cdot Y$ for vector fields on an open set of $\R^n$. The bracket is bilinear, antisymmetric, and satisfies the Jacobi identity
$[X,[Y,Z]]+[Y,[Z,X]]+[Z,[X,Y]]=0$.
\end{proposition}
\begin{proof}
At each point, $f\mapsto[X,Y]f(p)$ is linear, and the Leibniz rule holds because the terms $Xf\cdot Yg$ and $Yf\cdot Xg$ produced by expanding
$X(Y(fg))$ and $Y(X(fg))$ cancel. So $[X,Y]_p\in T_pM$. In a chart, $X(Yf)=\sum_{i,j}X^j\partial_j(Y^i\partial_if)=\sum_{i,j}X^j\partial_jY^i\,\partial_if+\sum_{i,j}X^jY^i\partial_j\partial_if$; the last
sum is symmetric in $X$ and $Y$ because mixed partial derivatives commute, so it cancels in the commutator, leaving the formula, whose coefficients are
smooth. Bilinearity and antisymmetry are clear, and the Jacobi identity is the identity $[A,[B,C]]+[B,[C,A]]+[C,[A,B]]=0$, true for any three operators
on any vector space, by expanding the twelve terms.
\end{proof}

The bracket measures how two flows fail to commute, and the cleanest way to see it is to run four flows in a loop.

\begin{theorem}[the bracket as a small loop]\label{thm:bracketloop}
Let $X,Y$ have flows $\varphi^X,\varphi^Y$ and let $p\in M$. For small $s$, let
$\psi(s)=\varphi^Y_{-s}\circ\varphi^X_{-s}\circ\varphi^Y_s\circ\varphi^X_s(p)$. Then for every smooth $f$,
\[
f(\psi(s))=f(p)+s^2\,[X,Y]_pf+O(s^3).
\]
In a chart around $p$: $\psi(s)=p+s^2[X,Y]_p+O(s^3)$.
\end{theorem}
\begin{proof}
Let $H(a,b,c,d)=f\bigl(\varphi^Y_d\circ\varphi^X_c\circ\varphi^Y_b\circ\varphi^X_a(p)\bigr)$, which is smooth near $0$ by Theorem~\ref{thm:flowsmooth}, so that
$f(\psi(s))=H(s,s,-s,-s)$. Call the variables $a,b,c,d$ the first to fourth, attached to the fields $V_1=X$, $V_2=Y$, $V_3=X$, $V_4=Y$. Setting all
variables but one to $0$ gives $H=f(\varphi^{V_i}_t(p))$, whose first two derivatives at $0$ are $V_if(p)$ and $V_i(V_if)(p)$, because
$\frac{\dd}{\dd t}g(\varphi_t(q))=(Vg)(\varphi_t(q))$ for any $g$. For a mixed derivative with $i<j$, set the other variables to $0$: then
$H=g_{t_j}(\varphi^{V_i}_{t_i}(p))$ with $g_{t_j}=f\circ\varphi^{V_j}_{t_j}$, and $\partial_{t_i}\partial_{t_j}H(0)=V_i(V_jf)(p)$: the field applied to the point
first acts on the function last. With $\varepsilon=(1,1,-1,-1)$, Taylor's formula gives $f(\psi(s))=f(p)+s\sum_i\varepsilon_iV_if(p)+\frac{s^2}2Q+O(s^3)$
with
\[
Q=\sum_iV_i^2f+2\sum_{i<j}\varepsilon_i\varepsilon_jV_iV_jf.
\]
The first-order term is $Xf+Yf-Xf-Yf=0$. In $Q$, the squares contribute $2X^2f+2Y^2f$, and the six mixed terms contribute
$2(XY-X^2-XY-YX-Y^2+XY)f$. The sum is $Q=2(XY-YX)f$, which is the claim.
\end{proof}

\begin{figure}[ht]
\centering
\fig{ch02-bracketloop}
\caption{The loop of Theorem~\ref{thm:bracketloop} for $X=\frac\partial{\partial x}$ and $Y=x\frac\partial{\partial y}$ in the plane, with $p=(0.6,0)$ and $s=0.6$.
Moving right, the vertical flow is stronger, so the upward step $s(x+s)$ exceeds the downward step $sx$ taken after moving back. The loop misses by
exactly $s^2[X,Y]=s^2\frac\partial{\partial y}$ (Exercise~\ref{ex:noncommute}).}
\label{fig:bracketloop}
\end{figure}

There is a third description, which explains the bracket's other name. A diffeomorphism $F$ pushes a vector field forward: $(F_*Y)_{F(p)}=dF_p(Y_p)$.

\begin{theorem}[Lie derivative]\label{thm:liederiv}
Let $X$ be complete with flow $\varphi_t$, and let $Y$ be a vector field. Then
\[
[X,Y]_p=\frac{\dd}{\dd t}\Big|_{t=0}\,d(\varphi_{-t})_{\varphi_t(p)}\bigl(Y_{\varphi_t(p)}\bigr),
\]
and more generally the derivative at time $t$ of the curve $t\mapsto d(\varphi_{-t})(Y_{\varphi_t(p)})$ in $T_pM$ is $d(\varphi_{-t})([X,Y]_{\varphi_t(p)})$.
\end{theorem}
\begin{proof}
Work in a chart around $p$, which contains $\varphi_t(p)$ and $\varphi_{-t}(q)$ for $q$ near $p$ and $t$ small. By Theorem~\ref{thm:flowsmooth} the map
$(s,q)\mapsto D_q\varphi(s,q)$ is smooth with value $I$ at $s=0$ and derivative $DX(q)$ in $s$ at $s=0$ (the variational equation), so
$D\varphi_{-t}(\varphi_t(p))=I-tDX(p)+O(t^2)$. Also $\varphi_t(p)=p+tX(p)+O(t^2)$, so $Y(\varphi_t(p))=Y(p)+tDY(p)X(p)+O(t^2)$. Multiplying,
\[
D\varphi_{-t}(\varphi_t(p))\,Y(\varphi_t(p))=Y(p)+t\bigl(DY\cdot X-DX\cdot Y\bigr)(p)+O(t^2),
\]
and Proposition~\ref{prop:bracket} identifies the coefficient of $t$. For the general time, the flow property gives
$d\varphi_{-t-h}(Y_{\varphi_{t+h}(p)})=d\varphi_{-t}\bigl(d\varphi_{-h}(Y_{\varphi_h(q)})\bigr)$ with $q=\varphi_t(p)$; differentiate in $h$ at $0$ and use the first part at $q$,
the map $d\varphi_{-t}$ being linear and independent of $h$.
\end{proof}

\begin{theorem}[commuting flows]\label{thm:commuting}
Let $X$ and $Y$ be complete. Their flows commute, $\varphi^X_t\circ\varphi^Y_s=\varphi^Y_s\circ\varphi^X_t$ for all $s,t$, if and only if $[X,Y]=0$.
\end{theorem}
\begin{proof}
If the flows commute, the loop $\psi(s)$ of Theorem~\ref{thm:bracketloop} is constant equal to $p$, so $[X,Y]_pf=0$ for all $f$ and $p$. Conversely, if
$[X,Y]=0$, Theorem~\ref{thm:liederiv} shows that $t\mapsto d\varphi^X_{-t}(Y_{\varphi^X_t(p)})$ has derivative zero, so it is constant, equal to $Y_p$: that is,
$d\varphi^X_t(Y_p)=Y_{\varphi^X_t(p)}$ for all $t$ and $p$. Fix $t$ and $q$, and let $c(s)=\varphi^X_t(\varphi^Y_s(q))$. Its velocity is
$d\varphi^X_t(Y_{\varphi^Y_s(q)})=Y_{c(s)}$, so $c$ is the integral curve of $Y$ starting at $\varphi^X_t(q)$, which is $\varphi^Y_s(\varphi^X_t(q))$.
\end{proof}

The same proofs give local versions: if $[X,Y]=0$ near $p$, then $\varphi^X_t\circ\varphi^Y_s=\varphi^Y_s\circ\varphi^X_t$ near $p$ for small $s,t$. That is the form in
which we use the theorem now.

\section{Frobenius' theorem}

A non-vanishing vector field can be straightened: near a point where $X_p\ne0$ there are coordinates in which $X=\frac\partial{\partial x^1}$ (this is the
case $k=1$ of the theorem below, and its proof is the same). For several vector fields, or more generally for a field of $k$-planes, there is an
obstruction, and it is the bracket.

\begin{definition}
A \emph{distribution} of rank $k$ on $M$ is a choice of a $k$-dimensional subspace $\Delta_p\subset T_pM$ at each point, which is smooth: near each point
there are $k$ vector fields that span $\Delta_q$ at every nearby $q$. It is \emph{involutive} if $[X,Y]$ takes values in $\Delta$ whenever $X$ and $Y$ do, and
\emph{integrable} if every point has a chart with coordinates $x^1,\dots,x^n$ in which $\Delta$ is spanned by $\frac\partial{\partial x^1},\dots,\frac\partial{\partial x^k}$.
\end{definition}

In the coordinates of an integrable distribution the slices $x^{k+1},\dots,x^n=\text{const}$ are $k$-dimensional submanifolds tangent to $\Delta$ at each of their
points: they are the \emph{integral manifolds}, and the distribution is the field of their tangent planes.

\begin{theorem}[Frobenius]\label{thm:frobenius}
A distribution is integrable if and only if it is involutive.
\end{theorem}
\begin{proof}
\emph{Integrable implies involutive.} In the coordinates of the definition, vector fields with values in $\Delta$ are the $\sum_{i\le k}a^i\partial_i$, and by the
formula of Proposition~\ref{prop:bracket} the bracket of two of them is $\sum_{i\le k}\bigl(a^j\partial_jb^i-b^j\partial_ja^i\bigr)\partial_i$, which is again in $\Delta$.

\emph{Involutive implies integrable.} The question is local, so we work in a chart around $p$ with $p=0$, on an open set of $\R^n$. After a linear change
of coordinates, $\Delta_0$ is complementary to $\{0\}\times\R^{n-k}$, so the projection $\pi\colon\R^n\to\R^k$ onto the first $k$ coordinates is an isomorphism
from $\Delta_0$ onto $\R^k$, and by continuity from $\Delta_q$ onto $\R^k$ for $q$ near $0$. (In terms of spanning fields $Z_1,\dots,Z_k$, the $k\times k$ matrix
of their first coordinates is invertible at $0$, hence near $0$.)

\emph{A special frame.} Let $X_i(q)$ be the unique vector of $\Delta_q$ with $\pi(X_i(q))=e_i$. Inverting the $k\times k$ matrix above expresses $X_i$ as a
combination of the $Z_j$ with smooth coefficients, so $X_i$ is a smooth field, of the form
\[
X_i=\partial_i+\sum_{l>k}a_i^l\,\partial_l.
\]
By the coordinate formula, $[X_i,X_j]$ has no component along $\partial_1,\dots,\partial_k$: those components are derivatives of the constants $0$ and $1$. By
involutivity $[X_i,X_j]$ lies in $\Delta$, and a vector of $\Delta_q$ with $\pi=0$ is zero. So the fields $X_1,\dots,X_k$ commute.

\emph{Coordinates from flows.} Define, for $t\in\R^k$ and $y\in\R^{n-k}$ near $0$,
\[
\Phi(t,y)=\varphi^{X_1}_{t_1}\circ\varphi^{X_2}_{t_2}\circ\dots\circ\varphi^{X_k}_{t_k}(0,y).
\]
Because the flows commute (in the local form of Theorem~\ref{thm:commuting}), any one of them can be brought to the front, so
$\partial_{t_i}\Phi(t,y)=X_i(\Phi(t,y))$ for every $i$. At the origin, $D\Phi(0)$ has columns $X_1(0),\dots,X_k(0)$, whose first $k$ coordinates form the
identity, and $\partial_y\Phi(0)=\bigl(\begin{smallmatrix}0\\ I\end{smallmatrix}\bigr)$. This block-triangular matrix is invertible, so by the inverse function
theorem $\Phi$ is a diffeomorphism near $0$, and $(t,y)=\Phi^{-1}$ is a chart. In this chart $\frac\partial{\partial t_i}=d\Phi(e_i)=X_i$, so $\Delta$ is spanned by
$\frac\partial{\partial t_1},\dots,\frac\partial{\partial t_k}$.
\end{proof}

\begin{beauty}
The theorem turns a question about solving a system of partial differential equations (does a family of surfaces with prescribed tangent planes exist?)
into an algebraic test on brackets, which one checks by differentiating. And the proof builds the surfaces by literally flowing along the given
directions: once the directions commute, running the flows for times $t_1,\dots,t_k$ in any order lands in the same place, so the times themselves
become coordinates.
\end{beauty}

\begin{example}[a non-integrable plane field]\label{ex:contact}
On $\R^3$ let $\Delta$ be spanned by $Y_1=\partial_y$ and $Y_2=\partial_x+y\,\partial_z$; it is the kernel of the form $\dd z-y\dd x$. Then $[Y_1,Y_2]=\partial_z\notin\Delta$,
so there is no surface tangent to $\Delta$, even locally. By Theorem~\ref{thm:bracketloop}, moving along $Y_1$ and $Y_2$ in a small loop produces a net
displacement in the forbidden direction $\partial_z$. This is how a car parks (Exercise~\ref{ex:parking}), and why a falling cat can turn over: the
controls span a non-integrable distribution, so small loops in the allowed directions reach every direction.
\end{example}

\section{The hairy ball theorem}

On the odd-dimensional sphere $S^{2m-1}\subset\R^{2m}$ there is a tangent vector field that vanishes nowhere:
$v(x_1,x_2,\dots,x_{2m-1},x_{2m})=(-x_2,x_1,\dots,-x_{2m},x_{2m-1})$, which is $x\mapsto ix$ in $\C^m$, the velocity of the rotation
$x\mapsto e^{it}x$. On even-dimensional spheres there is none. For the ordinary sphere this was known to Poincaré in 1885, through his theory of the
index of singular points; Brouwer proved it in every even dimension in 1912. The proof below is Milnor's, from 1978.

\begin{theorem}[hairy ball]\label{thm:hairyball}
If $n$ is even, every continuous tangent vector field on $S^n$ vanishes somewhere.
\end{theorem}
\begin{proof}
Write $N=n+1$, which is odd, and suppose $v$ is a continuous tangent field on $S^n$ that vanishes nowhere.

\emph{We may assume $v$ smooth and of length $1$.} Let $m=\min|v|>0$. By the Stone--Weierstrass theorem (Volume~7, Theorem~4.19) there is a polynomial map
$P\colon\R^N\to\R^N$ with $|P-v|<\frac m2$ on $S^n$. Its tangential part $w(x)=P(x)-\langle P(x),x\rangle x$ is smooth, tangent, and $|w-v|\le|P-v|<\frac m2$
because $v(x)\perp x$ and orthogonal projection onto $x^\perp$ does not increase distances. So $w$ vanishes nowhere, and $w/|w|$ is smooth, tangent and of
length $1$. Replace $v$ by it.

\emph{Extension.} On the open shell $A=\{\frac12<|x|<\frac32\}$ put $v(x)=|x|\,v(\frac x{|x|})$, a smooth field with $v(x)\perp x$ and $|v(x)|=|x|$. Let
$f_t(x)=x+tv(x)$. By Pythagoras $|f_t(x)|=\sqrt{1+t^2}\,|x|$, so $f_t$ maps the sphere of radius $r$ into the sphere of radius $r\sqrt{1+t^2}$.

\emph{$f_t$ is a diffeomorphism of $A$ onto the shell $\sqrt{1+t^2}\,A$ for small $t$.} Since $v$ is $C^1$ on a neighbourhood of the compact $\bar A$, it is
Lipschitz there with some constant $L$, so $|f_t(x)-f_t(y)|\ge(1-|t|L)|x-y|$ and $f_t$ is injective on $A$ for $|t|<\frac1L$. Its Jacobian $\det(I+tDv(x))$ is a
polynomial in $t$ with coefficients continuous in $x$, equal to $1$ at $t=0$; on the compact $\bar A$ it stays positive for $|t|$ small. For surjectivity,
consider $g_t(x)=\frac{f_t(x)}{\sqrt{1+t^2}}$ on the unit sphere: it maps $S^n$ into itself, and its differential at $x$ is the restriction of
$\frac{I+tDv(x)}{\sqrt{1+t^2}}$ to $T_xS^n$, injective for small $t$, hence an isomorphism onto $T_{g_t(x)}S^n$. By the inverse function theorem $g_t$ is a
local diffeomorphism, so $g_t(S^n)$ is open in $S^n$; it is also compact, hence closed; and $S^n$ is connected, so $g_t$ is onto. Since $f_t(rx)=rf_t(x)$, the
image $f_t(A)$ is the whole shell $\sqrt{1+t^2}A$. An injective smooth map with invertible differential is a diffeomorphism onto its open image, by the
inverse function theorem.

\emph{Two computations of a volume.} By the change of variables formula (Volume~7, Theorem~9.8),
\[
\lambda\bigl(f_t(A)\bigr)=\int_A\det\bigl(I+tDv(x)\bigr)\dd x,
\]
which is a polynomial in $t$, because the determinant is. On the other hand $f_t(A)=\sqrt{1+t^2}\,A$ has volume $(1+t^2)^{N/2}\lambda(A)$. So
$(1+t^2)^{N/2}$ coincides with a polynomial $R(t)$ for all small $t$. Then $R(t)^2=(1+t^2)^N$ on an interval, hence identically, as polynomials. But
$(1+t^2)^N=(t-i)^N(t+i)^N$ has the root $i$ with odd multiplicity $N$, while every root of $R^2$ has even multiplicity. This contradiction proves the theorem.
\end{proof}

\begin{beauty}
The proof never mentions topology. It moves every point of a shell a little along the hair, notices that the shell is carried onto a slightly larger
shell, and compares two formulas for the new volume: one is a polynomial in $t$, the other is $(1+t^2)^{N/2}$, which is a polynomial exactly when $N$ is
even. The parity of the dimension, which is all the theorem is about, enters through the square root.
\end{beauty}

\begin{remark}
On the Earth, a continuous wind field has a calm point somewhere. How many independent nowhere-vanishing fields an odd sphere $S^{n-1}$ carries was
settled by J.~F.~Adams in 1962: if $n=2^{4a+b}\cdot\text{odd}$ with $0\le b\le3$, the maximum is $2^b+8a-1$. In particular only $S^1$, $S^3$ and $S^7$ have
$n-1$ independent fields, a fact tied to the existence of the complex numbers, the quaternions and the octonions (Exercise~\ref{ex:hairyodd}).
Chapter~6 will reprove the hairy ball theorem with homology.
\end{remark}

\section{Exercises}

\begin{exercise}[$\star$]\label{ex:noncomplete}
\emph{A field that is not complete.} (a) Find the flow of $X=x^2\frac{\dd}{\dd x}$ on $\R$ and its flow domain. (b) Show that the field $x^2\frac{\dd}{\dd x}$
extends to a smooth field on the circle $\R P^1=\R\cup\{\infty\}$, which is complete, and describe its flow.
\end{exercise}

\begin{exercise}[$\star$]\label{ex:sqrtuniqueness}
\emph{Uniqueness needs Lipschitz.} Show that for every $c\ge0$ the function $x_c(t)=\frac{(t-c)^2}4$ for $t\ge c$, $0$ for $t\le c$, solves $x'=\sqrt{|x|}$ with
$x(0)=0$. Why does this not contradict Theorem~\ref{thm:picard}? Show on the other hand that $x'=x\log|x|$ (extended by $0$ at $0$), which is not Lipschitz
at $0$ either, has unique solutions. (Separate variables.)
\end{exercise}

\begin{exercise}[$\star$]\label{ex:noncommute}
\emph{Flows that do not commute.} Let $X=\partial_x$ and $Y=x\,\partial_y$ on $\R^2$. (a) Compute their flows and $[X,Y]$. (b) Compute the loop $\psi(s)$ of
Theorem~\ref{thm:bracketloop} exactly and compare with the theorem (Figure~\ref{fig:bracketloop}). (c) Show that $\Delta=\operatorname{span}(X,Y)$ is not a
distribution of rank $2$ on $\R^2$ and that Frobenius' theorem does not apply. Where does it fail?
\end{exercise}

\begin{exercise}[$\star$]\label{ex:linearfields}
\emph{Linear vector fields.} For $A\in M_n(\R)$ let $X_A(x)=Ax$. (a) Show that the flow of $X_A$ is $x\mapsto e^{tA}x$. (b) Show that $[X_A,X_B]=X_{BA-AB}$.
(c) Deduce that $e^{tA}$ and $e^{sB}$ commute for all $s,t$ if and only if $AB=BA$. (d) For $A=\bigl(\begin{smallmatrix}0&-1\\1&0\end{smallmatrix}\bigr)$, $B=\bigl(\begin{smallmatrix}1&0\\0&-1\end{smallmatrix}\bigr)$ compute
$e^{-sB}e^{-sA}e^{sB}e^{sA}$ to second order and check Theorem~\ref{thm:bracketloop}.
\end{exercise}

\begin{exercise}[$\star\star$]\label{ex:parking}
\emph{Parallel parking.} A unicycle in the plane has configuration $(x,y,\theta)\in\R^2\times S^1$: position and heading. It can drive,
$F=\cos\theta\,\partial_x+\sin\theta\,\partial_y$, and turn, $T=\partial_\theta$, but not slide sideways. (a) Compute $[T,F]$ and interpret it. (b) Starting
from $(0,0,0)$, turn by $s$, drive forward $s$, turn back by $s$, drive backward $s$. Compute the exact final configuration and compare with
Theorem~\ref{thm:bracketloop}. (c) Show that the distribution spanned by $F$ and $T$ is not integrable, and deduce that no function $g(x,y,\theta)$ other than
the constants is preserved by all motions. (This is why a car can be parked in a space barely longer than itself.)
\end{exercise}

\begin{exercise}[$\star$]\label{ex:hairyodd}
\emph{Odd spheres.} (a) Check that $v(x)=ix$ is a unit tangent field on $S^{2m-1}\subset\C^m$. (b) On $S^3\subset\mathbb{H}$, the quaternions of Volume~3, show that
$x\mapsto xi$, $x\mapsto xj$, $x\mapsto xk$ are three tangent fields, orthonormal at every point. (c) Show that the flows of the three fields in (b) are
$x\mapsto xe^{ti}$, $x\mapsto xe^{tj}$, $x\mapsto xe^{tk}$, and that the fields $x\mapsto ix$ (left multiplication) commute with all three.
\end{exercise}

\begin{exercise}[$\star\star$]\label{ex:straighten}
\emph{Straightening a field.} Let $X$ be a vector field with $X_p\ne0$. Prove directly, following the proof of Theorem~\ref{thm:frobenius} with $k=1$, that
there is a chart around $p$ in which $X=\frac\partial{\partial x^1}$. Find such coordinates for the rotation field $X=-y\,\partial_x+x\,\partial_y$ near $(1,0)$.
\end{exercise}

\section{Hints}

\begin{hint}{ex:noncomplete}
(a) Separate variables. (b) Use the chart $u=-\frac1x$ near $\infty$.
\end{hint}
\begin{hint}{ex:sqrtuniqueness}
$\sqrt{|x|}$ has infinite slope at $0$. For $x\log|x|$, put $u=\log|x|$.
\end{hint}
\begin{hint}{ex:noncommute}
$\varphi^Y_s(x,y)=(x,y+sx)$. For (c), look at the points where $Y=0$.
\end{hint}
\begin{hint}{ex:linearfields}
For (b) use $[X,Y]=DY\cdot X-DX\cdot Y$ and $DX_A=A$. For (d), $e^{sA}=I+sA+\frac{s^2}2A^2+O(s^3)$.
\end{hint}
\begin{hint}{ex:parking}
Turning does not move the position. For (c), a preserved function is killed by $F$ and $T$, hence by $[T,F]$.
\end{hint}
\begin{hint}{ex:hairyodd}
A vector $w$ is tangent to the unit sphere at $x$ if $\operatorname{Re}\langle x,w\rangle=0$; for quaternions, $\langle x,xq\rangle=|x|^2\langle1,q\rangle$.
\end{hint}
\begin{hint}{ex:straighten}
Take $\Phi(t,y)=\varphi^X_t(0,y)$ in coordinates where $X_p=e_1$. For the rotation, polar coordinates.
\end{hint}

\section{Solutions}

\begin{solution}{ex:noncomplete}
\idea Solutions of $x'=x^2$ reach infinity in finite time; on the circle, infinity is just another point, and the solution goes through it.
\writeup (a) For $x\ne0$, $(\frac1x)'=-1$, so $\frac1{x(t)}=\frac1{x_0}-t$ and $\varphi(t,x_0)=\frac{x_0}{1-tx_0}$, a formula that also covers $x_0=0$. The solution
blows up at $t=\frac1{x_0}$, so $J(x_0)=\left]-\infty,\frac1{x_0}\right[$ for $x_0>0$, $\left]\frac1{x_0},\infty\right[$ for $x_0<0$, and $\R$ for $x_0=0$. The flow domain is
$\{(t,x):tx<1\}$, open as Theorem~\ref{thm:flowcont} says, and not all of $\R\times\R$. (b) In the chart $u=-\frac1x$ around $\infty$ (with $\infty\mapsto0$), the
chain rule gives $u'=\frac{x'}{x^2}=1$: the field is $\frac{\dd}{\dd u}$, smooth and non-zero near $\infty$. On the compact manifold $\R P^1$ it is complete
(Theorem~\ref{thm:complete}): its flow is the Möbius map $x\mapsto\frac x{1-tx}$, which sends $\frac1t$ to $\infty$ and $\infty$ to $-\frac1t$. A solution that ``explodes''
on the line just passes through the point at infinity and comes back from the other side. The only fixed point is $0$.
\end{solution}

\begin{solution}{ex:sqrtuniqueness}
\idea Non-uniqueness needs a solution that can leave the equilibrium in finite time, which means $\int_0\frac{\dd x}{F(x)}<\infty$.
\writeup For $t>c$, $x_c'=\frac{t-c}2=\sqrt{x_c}$; for $t<c$ both sides vanish; and at $t=c$ both one-sided derivatives are $0$. So each $x_c$ is a solution, and
so is $0$: infinitely many. Theorem~\ref{thm:picard} needs a Lipschitz bound near $0$, and $\frac{\sqrt{|x|}}{|x|}\to\infty$. The time to leave $0$ is governed by
$\int_0^\varepsilon\frac{\dd x}{\sqrt x}<\infty$. For $F(x)=x\log|x|$, the integral $\int_0^\varepsilon\frac{\dd x}{x\log x}=[\log|\log x|]$ diverges, so no solution reaches or leaves $0$
in finite time. Concretely, for $x>0$ put $u=\log x$: then $u'=\frac{x'}x=u$, so $u=u_0e^t$ and $x(t)=\exp(e^t\log x_0)$, defined for all $t$ and never $0$; a solution
that is $0$ at one time is $0$ at all times, since it cannot come from or go to a positive value. So uniqueness holds although $F$ is not Lipschitz:
the Lipschitz condition is sufficient, not necessary (the sharp condition is Osgood's, $\int_0\frac{\dd r}{\omega(r)}=\infty$ for the modulus of continuity $\omega$).
\end{solution}

\begin{solution}{ex:noncommute}
\idea Each flow is explicit, so the loop can be computed exactly; the vertical flow speeds up as one moves right, which is what the bracket detects.
\writeup (a) $\varphi^X_t(x,y)=(x+t,y)$ and $\varphi^Y_s(x,y)=(x,y+sx)$. By the formula, $[X,Y]=DY\cdot X-DX\cdot Y$ with $Y(x,y)=(0,x)$, $DY\cdot X=\partial_xY=(0,1)$ and
$DX=0$, so $[X,Y]=\partial_y$. (b) From $p=(x,y)$: $(x+s,y)$, then $(x+s,y+s(x+s))$, then $(x,y+s(x+s))$, then $(x,y+s(x+s)-sx)=(x,y+s^2)$. So
$\psi(s)=p+s^2\partial_y$ exactly, with no $O(s^3)$ term. In particular $\varphi^X$ and $\varphi^Y$ do not commute. (c) At points with $x=0$, $Y=0$ and $\Delta$ is
spanned by $X$ alone: the rank drops, so $\Delta$ is not a distribution of rank $2$. Away from $x=0$, $X$ and $Y$ span the whole plane, an integrable
distribution of rank $2$ (the plane is its own integral manifold). Along the line $x=0$, the loop leaves the line in the direction $\partial_y$, which
is tangent to the line: nothing contradicts Frobenius, which assumes the rank is constant.
\end{solution}

\begin{solution}{ex:linearfields}
\idea Linear fields are the infinitesimal version of matrices, and the bracket of fields becomes a commutator of matrices, with a sign that comes from the
order in which compositions act on points.
\writeup (a) $\frac{\dd}{\dd t}e^{tA}x=Ae^{tA}x$ and $e^{0}x=x$, so by uniqueness this is the flow (Volume~3, Chapter~3, for the exponential). (b) $DX_B=B$, so
$[X_A,X_B](x)=B(Ax)-A(Bx)=(BA-AB)x$. (c) By Theorem~\ref{thm:commuting}, the flows commute iff $[X_A,X_B]=0$ iff $AB=BA$. (d) The loop of Theorem~\ref{thm:bracketloop} with $X=X_A$ (applied first) and $Y=X_B$ is $x\mapsto e^{-sB}e^{-sA}e^{sB}e^{sA}x$. Write the four
factors, from left to right, as $e^{sM_i}$ with $(M_1,M_2,M_3,M_4)=(-B,-A,B,A)$ and expand each as $I+sM_i+\frac{s^2}2M_i^2+O(s^3)$. The first-order term
is $\sum M_i=0$. The second-order term is $\frac12\sum_iM_i^2+\sum_{i<j}M_iM_j=(A^2+B^2)+(BA-B^2-BA-AB-A^2+BA)=BA-AB$. So the loop is
$I+s^2(BA-AB)+O(s^3)$, as the theorem predicts from $[X_A,X_B]=X_{BA-AB}$. For the given matrices, $AB=\bigl(\begin{smallmatrix}0&1\\1&0\end{smallmatrix}\bigr)$ and
$BA=\bigl(\begin{smallmatrix}0&-1\\-1&0\end{smallmatrix}\bigr)$, so the loop is $I-2s^2\bigl(\begin{smallmatrix}0&1\\1&0\end{smallmatrix}\bigr)+O(s^3)$: a rotation followed by a squeeze,
undone in the opposite order, leaves a small shear. The sign is worth noticing: $X_A\mapsto A$ reverses brackets. Chapter~7 meets the same sign in the
Lie algebra of a matrix group, where the natural fields are $x\mapsto xA$, and then the sign disappears.
\end{solution}

\begin{solution}{ex:parking}
\idea The bracket of turning and driving is a sideways slide, a direction the unicycle cannot take directly but can reach by a small loop.
\writeup (a) In the coordinates $(x,y,\theta)$ we have $F=(\cos\theta,\sin\theta,0)$, $T=(0,0,1)$, $DT=0$ and $DF\cdot T=\partial_\theta F=(-\sin\theta,\cos\theta,0)$, so
$[T,F]=-\sin\theta\,\partial_x+\cos\theta\,\partial_y$: the unit vector perpendicular to the heading, to the left. (b) Turning by $s$ gives $(0,0,s)$; driving $s$ gives
$(s\cos s,s\sin s,s)$; turning back gives $(s\cos s,s\sin s,0)$; reversing $s$ gives \[(s\cos s-s,\ s\sin s,\ 0)=\bigl(-\tfrac{s^3}2,\ s^2,\ 0\bigr)+O(s^4).\] Theorem~\ref{thm:bracketloop}
with $X=T$ (applied first) and $Y=F$ predicts $s^2[T,F]=s^2\partial_y$ at $\theta=0$: a sideways displacement $s^2$ to the left, with an error $O(s^3)$, here
$-\frac{s^3}2$ in $x$. (c) $F$, $T$ and $[T,F]$ are independent at every point, so $[T,F]\notin\operatorname{span}(F,T)$: the distribution is not involutive, hence not
integrable. If $g$ were preserved by both motions, then $Fg=Tg=0$, hence $[T,F]g=T(Fg)-F(Tg)=0$, so $dg$ vanishes on three independent vectors at each point:
$dg=0$ and $g$ is locally constant. A car (with one more degree of freedom, the steering angle) behaves the same way, and parallel parking is a sequence
of such loops, each gaining a sideways displacement of the order of the square of its size.
\end{solution}

\begin{solution}{ex:hairyodd}
\idea On spheres sitting in a number system, multiplication by a unit imaginary number is a rotation that moves every point perpendicular to itself.
\writeup (a) $|ix|=|x|=1$ and $\operatorname{Re}\langle x,ix\rangle=\operatorname{Re}\sum_j\bar x_jix_j=\operatorname{Re}(i|x|^2)=0$, so $ix\perp x$ in the real inner product. (b) Left multiplication
by $x$ multiplies all lengths by $|x|$ (as $|xa|=|x||a|$), so $\langle xa,xb\rangle=|x|^2\langle a,b\rangle$ by polarization. With $a,b\in\{1,i,j,k\}$
orthonormal and $|x|=1$, the vectors $x,xi,xj,xk$ are orthonormal: the last three are tangent at $x$ and orthonormal. So $S^3$ carries three everywhere
independent fields, and its tangent spaces are all identified with $\R^3$ at once. (c) $\frac{\dd}{\dd t}xe^{ti}=xe^{ti}i$, which is the field $y\mapsto yi$
at $y=xe^{ti}$; so this is the flow. Left and right multiplications commute by associativity, $(ix)q=i(xq)$, so the flows $x\mapsto e^{ti}x$ and $x\mapsto xe^{sq}$ commute,
and by Theorem~\ref{thm:commuting} the brackets of the left field with the right fields vanish. The right fields among themselves do not commute: their
brackets are again right fields, as Chapter~7 will compute ($\mathfrak{su}(2)$).
\end{solution}

\begin{solution}{ex:straighten}
\idea Flowing from a hypersurface transverse to $X$ turns the time of flow into a coordinate along which $X$ is the derivative.
\writeup Choose a chart with $p=0$ and $X_0=e_1$, and let $\Phi(t,y)=\varphi^X_t(0,y)$ for $t\in\R$ and $y\in\R^{n-1}$ small. Then $\partial_t\Phi=X(\Phi)$, and
$D\Phi(0)=(X_0\mid\begin{smallmatrix}0\\ I\end{smallmatrix})=I$. By the inverse function theorem $\Phi$ is a local diffeomorphism and in the chart $\Phi^{-1}$ we have
$\frac\partial{\partial t}=X$. No bracket condition is needed when $k=1$, since $[X,X]=0$. For the rotation field near $(1,0)$, the flow is the rotation by $t$, and
$\Phi(t,y)=((1+y)\cos t,(1+y)\sin t)$ (flowing from the horizontal line through $(1,0)$). The new coordinates are $t=\arg$ and $y=r-1$: polar coordinates,
in which $-y\partial_x+x\partial_y=\partial_\theta$.
\end{solution}

\section*{Sources and further reading}
\addcontentsline{toc}{section}{Sources and further reading}

The account of Lie's arrest and of Darboux's intervention is from A.~Stubhaug, \emph{The Mathematician Sophus Lie: It Was the Audacity of My Thinking}
(Springer, 2002), Chapter~13, and from the MacTutor History of Mathematics archive, \emph{Marius Sophus Lie}; Lie's remark about mathematicians and
prison comes from his letters as quoted by Stubhaug, and we have paraphrased it. T.~Hawkins, \emph{Emergence of the Theory of Lie Groups} (Springer, 2000), traces how Lie's
integration theory of differential equations grew into the theory of groups.

\'E.~Picard, ``M\'emoire sur la th\'eorie des \'equations aux d\'eriv\'ees partielles et la m\'ethode des approximations successives'', \emph{J.\ Math.\
Pures Appl.} (4) 6 (1890), 145--210; E.~Lindel\"of, ``Sur l'application de la m\'ethode des approximations successives aux \'equations diff\'erentielles
ordinaires du premier ordre'', \emph{C.\ R.\ Acad.\ Sci.\ Paris} 118 (1894), 454--457; T.~H.~Gr\"onwall, ``Note on the derivatives with respect to a
parameter of the solutions of a system of differential equations'', \emph{Ann.\ of Math.} 20 (1919), 292--296, where the inequality appears in the service
of exactly our Theorem~\ref{thm:flowsmooth}.

G.~Frobenius, ``\"Uber das Pfaffsche Problem'', \emph{J.\ reine angew.\ Math.} 82 (1877), 230--315; the theorem had been proved earlier by F.~Deahna,
\emph{J.\ reine angew.\ Math.} 20 (1840), 340--349, and used by Clebsch; Hawkins tells the story. J.~Milnor, ``Analytic proofs of the `hairy ball theorem'
and the Brouwer fixed point theorem'', \emph{Amer.\ Math.\ Monthly} 85 (1978), 521--524. J.~F.~Adams, ``Vector fields on spheres'', \emph{Ann.\ of Math.}
75 (1962), 603--632.

For the theory, J.~M.~Lee, \emph{Introduction to Smooth Manifolds}, 2nd ed. (Springer, 2013), Chapters~8, 9 and 19; for the differential equations,
M.~W.~Hirsch, S.~Smale and R.~L.~Devaney, \emph{Differential Equations, Dynamical Systems, and an Introduction to Chaos}, 3rd ed. (Academic Press, 2013),
Chapter~17, whose proof of smooth dependence we followed. The car as a control system with Lie brackets is in E.~Nelson, \emph{Tensor Analysis}
(Princeton, 1967), and the general theory of reachability by brackets (the theorems of Chow and of Rashevskii) in R.~Montgomery, \emph{A Tour of
Subriemannian Geometries} (AMS, 2002).
